% The ICC step by step, with a made-up example of four people (numbers checked against medis_compare.icc,
% 2026-11-20: ICC(C,1) 0.899, ICC(A,1) 0.863). Formulas from Koo and Li 2016, Table 3 (papers/).
\documentclass[11pt]{article}
% Week 12 preamble (copied 2026-11-19 from outdated/thesis_retired/week10/setup/preamble.tex, paths for
% weekly/week12/). Figures resolve from week12/week12_figures/.
\usepackage[T1]{fontenc}
\usepackage[utf8]{inputenc}
\usepackage{lmodern}
\usepackage[english]{babel}
\usepackage{amsmath,amssymb,mathtools}
\usepackage{graphicx}
\usepackage{float}
\usepackage[section]{placeins}  % every figure of a section is shown before the next section starts
\usepackage{booktabs,tabularx}
\usepackage{xcolor}
\usepackage{enumitem}
\usepackage[margin=2.5cm]{geometry}
\usepackage[hidelinks]{hyperref}
\usepackage{cleveref}
\graphicspath{{./}{week12_figures/}}
\providecommand{\clarify}[1]{\textcolor{red}{\textbf{Clarification needed:} #1}}
\providecommand{\todonote}[1]{\textcolor{blue}{\textbf{To do:} #1}}
\providecommand{\pending}[1]{\textcolor{blue}{\textbf{[#1]}}}  % a number waiting for the rerun of medis_compare.py

\setlength{\parindent}{0pt}
\setlength{\parskip}{6pt}

\begin{document}

\section*{The ICC step by step}

\subsection*{Step 0: the idea}

Every measurement is a value that belongs to the person plus an error of the method:
\[
  \text{tracer}_i = T_i + e^{\text{tracer}}_i, \qquad \text{Medis}_i = T_i + c + e^{\text{Medis}}_i .
\]
$T_i$ is what the two methods have in common for person $i$, $c$ is a constant bias, and both methods have their
own error. Nobody knows $T_i$. The ICC only asks how large the spread of $T$ between people is compared with
the errors (Koo and Li 2016, Introduction):
\[
  \text{ICC} = \frac{\text{variance between people}}{\text{variance between people} + \text{error variance}}.
\]

\subsection*{The example}

Four people, LM angle in degrees:

\begin{center}
\begin{tabular}{lrrrr}
  \toprule
  Person & Tracer & Medis & Pair mean $m_i$ & Difference $d_i$ = tracer $-$ Medis \\
  \midrule
  1 & 62 & 58 & 60 & 4 \\
  2 & 80 & 70 & 75 & 10 \\
  3 & 66 & 68 & 67 & $-2$ \\
  4 & 52 & 48 & 50 & 4 \\
  \midrule
  Mean & 65 & 61 & 63 & 4 \\
  \bottomrule
\end{tabular}
\end{center}

The two columns say everything:
\begin{itemize}
  \item the \textbf{pair means} $m_i$ show how the people differ (person 2 has a wide angle, person 4 a narrow one);
  \item the \textbf{differences} $d_i$ show where the methods disagree.
\end{itemize}

\subsection*{Step 1: the bias}

The mean difference is $\bar d = 4^\circ$: the tracer is on average $4^\circ$ higher. This is the constant $c$.
The consistency ICC does not count it (a shift of every angle by $4^\circ$ does not change who has the wider angle).

\subsection*{Step 2: the error variance}

Take the differences around their mean, so the bias is removed: $4-4=0$, $10-4=6$, $-2-4=-6$, $4-4=0$.
\[
  \operatorname{var}(d) = \frac{0^2 + 6^2 + (-6)^2 + 0^2}{4-1} = \frac{72}{3} = 24 .
\]
A difference holds the error of \emph{two} measurements, $e^{\text{tracer}} - e^{\text{Medis}}$. If the two
errors are independent, their variances add, so one measurement has on average half of it:
\[
  \sigma^2_{\text{error}} = \frac{24}{2} = 12 \qquad (\text{this is } MS_E).
\]
Both methods may be wrong; $\sigma^2_{\text{error}}$ is the \emph{average} error variance of the two.

\subsection*{Step 3: the variance between people}

Take the pair means around their mean 63: $-3,\ 12,\ 4,\ -13$.
\[
  \operatorname{var}(m) = \frac{9 + 144 + 16 + 169}{4-1} = \frac{338}{3} = 112.7 .
\]
A pair mean still carries some error: the mean of two measurements has error variance
$\sigma^2_{\text{error}}/2 = 6$. Remove it:
\[
  \sigma^2_{\text{between}} = 112.7 - 6 = 106.7 .
\]

\subsection*{Step 4: the ICC}

\[
  \text{ICC} = \frac{\sigma^2_{\text{between}}}{\sigma^2_{\text{between}} + \sigma^2_{\text{error}}}
             = \frac{106.7}{106.7 + 12} = 0.90 .
\]
90\% of the variance of one measurement is the people differing; 10\% is the methods disagreeing.

\subsection*{The same, in the notation of the paper}

Koo and Li (2016, Table 3) write the two-way mixed, consistency, single-measurement ICC, ICC(3,1), with mean
squares from an analysis of variance (rows = people, columns = tracer and Medis, $k = 2$ methods, $n = 4$ people):
\[
  \text{ICC}(3,1) = \frac{MS_R - MS_E}{MS_R + (k-1)\,MS_E}.
\]
With two methods the mean squares are the variances above:
\[
  MS_R = 2 \operatorname{var}(m) = 225.3, \qquad MS_E = \tfrac12 \operatorname{var}(d) = 12,
\]
\[
  \text{ICC}(3,1) = \frac{225.3 - 12}{225.3 + 12} = \frac{213.3}{237.3} = 0.90 .
\]
It is what \texttt{icc()} in \texttt{medis\_compare.py} computes.

\subsection*{Two routes, one answer}

Both routes remove the error that is still inside the pair means. They differ only in the scale they work on.

\begin{center}
\begin{tabular}{lll}
  \toprule
  & Example (Steps 3 and 4) & Paper (Table 3) \\
  \midrule
  Scale & one pair mean (half the error) & one measurement (full error) \\
  Pair means & $\operatorname{var}(m) = 112.7$ & $MS_R = 2\operatorname{var}(m) = 225.3$ \\
  Remove the error & $112.7 - \tfrac{12}{2} = 106.7$ & $MS_R - MS_E = 225.3 - 12 = 213.3$ \\
  Denominator & $106.7 + 12 = 118.7$ & $MS_R + MS_E = 225.3 + 12 = 237.3$ \\
  ICC & $106.7 / 118.7 = 0.90$ & $213.3 / 237.3 = 0.90$ \\
  \bottomrule
\end{tabular}
\end{center}

Every number of the paper route is twice the number of the example route ($213.3 = 2 \times 106.7$,
$237.3 = 2 \times 118.7$). The 2 is in both the numerator and the denominator, so it cancels:
\[
  \frac{MS_R - MS_E}{MS_R + MS_E}
  = \frac{2\,\sigma^2_{\text{between}}}{2\,\sigma^2_{\text{between}} + 2\,\sigma^2_{\text{error}}}
  = \frac{\sigma^2_{\text{between}}}{\sigma^2_{\text{between}} + \sigma^2_{\text{error}}} .
\]
The paper multiplies by $k$ (here 2) because each person has $k$ values in the table and the pair mean holds only
$1/k$ of the error. Multiplying by $k$ puts the pair means on the scale of one measurement, so a whole $MS_E$ can
be subtracted. This gives one formula for any number of methods; the example avoids the factor because
``subtract half the error'' is easier to follow with two.

\subsection*{If the bias is counted: absolute agreement}

The bias enters through $MS_C = n \cdot \sum (\text{column mean} - \text{grand mean})^2 = 4\,(2^2 + 2^2) = 32$:
\[
  \text{ICC}(A,1) = \frac{MS_R - MS_E}{MS_R + (k-1)MS_E + \frac{k}{n}(MS_C - MS_E)}
                  = \frac{213.3}{237.3 + \frac{2}{4}(32-12)} = \frac{213.3}{247.3} = 0.86 .
\]
Lower than 0.90, because the $4^\circ$ bias now counts as disagreement.

\subsection*{Why this form: two-way mixed effects, consistency, single measurement}

McGraw and Wong define 10 forms of the ICC. Koo and Li (2016, section ``How to Select the Correct ICC Form for
Interrater Reliability Studies'') choose between them with questions about the design. Each part of the name
answers one.

\textbf{Two-way (the model).} The table varies in two directions: rows (people) and columns (raters), and the
analysis of variance estimates both, how much people differ ($MS_R$) and how much the raters differ on average
($MS_C$). In a one-way model each person is measured by different raters (the paper's example is a multicentre
study with its own raters per centre), so the columns mean nothing and only the rows are modelled. Here every
artery pair is measured by the same two methods, the tracer and Medis, so the column effect (the bias) can be
estimated: two-way.

\textbf{Mixed effects (random or fixed raters).} People are always a random sample. The question is the raters.
In the random model they are a sample from many possible raters (two radiologists standing for all radiologists).
The paper uses the mixed model ``if the selected raters are the only raters of interest. With this model, the
results only represent the reliability of the specific raters involved''. Only the tracer and Medis are of
interest, not centerline software in general: mixed. The random and the mixed model use the same formula; the
difference ``is not on the calculation but on the experimental design of the reliability study and the
interpretation of the results''. Mixed changes the claim, not the number.

\textbf{Consistency (does a constant bias count).} ``Consistency definition concerns the degree to which one
rater's score (y) can be equated to another rater's score (x) plus a systematic error (c) (ie, y = x + c), whereas
absolute agreement concerns about the extent to which y equals x.'' The model uses geometry per SD, and adding a
constant to every value changes neither the SD nor anyone's place, so the hazard ratio is unchanged. A bias that
does not harm the model should not lower the reliability: consistency. The bias is reported on its own (the
Bland--Altman bias).

\textbf{Single measurement (one method or the mean of both).} ``If we plan to use the measurement from a single
rater as the basis of the actual measurement, `single rater' type should be selected even though the reliability
experiment involves 2 or more raters.'' The cohort gets one value per person, the tracer's; Medis does not exist
for the plaque-free cohort, so there is no average. The mean-of-two ICC is always higher, because averaging cancels
part of the error, and would overstate the reliability of the values the model uses: single.

\begin{center}
\small
\begin{tabular}{lll}
  \toprule
  Question & Answer here & Part of the name \\
  \midrule
  The same raters for all people? & Yes, tracer and Medis & Two-way \\
  Raters a sample, or the only ones of interest? & The only ones of interest & Mixed effects \\
  Does a constant bias matter? & No, the model uses per-SD values & Consistency \\
  One measurement or the mean of the raters? & One, the tracer value & Single \\
  \bottomrule
\end{tabular}
\end{center}

Together this is ICC(3,1) in the notation of Shrout and Fleiss. Consistency is the choice that needs defending:
the paper says absolute agreement ``should always be chosen'' for test--retest and intrarater studies. This
comparison is neither (two methods, not one method repeated), and the reason for consistency is the per-SD hazard
ratio; the absolute-agreement ICC is computed as well and is only slightly lower in the example (0.86 against 0.90).

\subsection*{What the ICC cannot tell}

\begin{itemize}
  \item \textbf{Which method is wrong.} $\sigma^2_{\text{error}}$ is the average of the two methods' errors. If
    Medis is the noisier one, the tracer alone is more reliable than the ICC says, and the other way round.
  \item \textbf{Errors both methods share.} If both stop early in the same faint vessel, the error is the same in
    both columns, so it ends up in the pair mean, not in the difference. It is then counted as people differing,
    and the ICC is too high.
  \item \textbf{Other samples.} With the same error, people who differ more give a higher ICC. These ICCs come
    from participants with plaque.
\end{itemize}

\end{document}
